%
% Slide template for the FreeBSD Developer Summit.
% Take it or leave it :-)
%
% It requires LaTeX and LaTeX Beamer [1] to compile.
% pdfLaTeX is recommended for compilation as it produces
% PDF file immediately.
%
% $ pdflatex some.latex
%
%
% It is also recommended to convert images to PDF
% by using ImageMagick ("convert") before including them.
%
% [1] http://latex-beamer.sourceforge.net
%

\documentclass{beamer}

\usepackage{url}
\usepackage[utf8]{inputenc}
\usepackage[russian]{babel}
\usepackage{verbatim}
\usepackage{graphicx}
\usepackage{listings}
\usepackage{mathtools}
\usepackage{color}

\mode<presentation>
{
  \definecolor{beamer@gker}{rgb}{0.8,0.0,0.0}
  \setbeamercolor*{structure}{fg=beamer@gker}
  \logo{\includegraphics[scale=0.5]{logo.pdf}}
}

\setbeamertemplate{footline}[text line]{%
  \parbox{\linewidth}{\vspace*{-8pt}
  \url{http://highsecure.ru/pkg}
  \hfill
  \insertshortauthor
  \hfill
  \insertframenumber\ of \inserttotalframenumber}}
\setbeamertemplate{navigation symbols}{}

\definecolor{light-gray}{gray}{0.60}

\title{Pkg - система управления пакетами FreeBSD}
\author{Всеволод Стахов \\ \url{vsevolod@FreeBSD.org}}
\institute{\includegraphics[scale=0.5]{logo.pdf}}
\date{ruBSD conference 14 Декабря 2013}

\begin{document}

\begin{frame}[plain]
  \titlepage
\end{frame}

\begin{frame}
\frametitle{Порты и пакеты}
Порты - основа для создания пакетов.
\begin{itemize}
  \item Устоявшаяся система
  \item Понятные и четкие правила
  \item Простые в создании и управлении (не всегда)
  \item Множество настроек
\end{itemize}
\end{frame}

\begin{frame}[fragile]
\frametitle{Порты без pkg}
\begin{figure}[h!]
  \centering
  \includegraphics[width=0.95\textwidth]{q1.eps}
\end{figure}
\end{frame}

\begin{frame}
\frametitle{Недостатки предыдущей схемы}
\begin{itemize}
  \item Make не может полноценно разрешать зависимости и конфликты
  \item Сложно обновлять и держать в актуальном состоянии
  \item Проблемы при миграции между версиями ОС
  \item Процесс компиляции из исходников занимает время
\end{itemize}
\end{frame}

\begin{frame}
\frametitle{Ожидаемая архитектура взаимодействия pkg и портов}
\begin{figure}[h!]
  \centering
  \includegraphics[width=0.99\textwidth]{q2.eps}
\end{figure}
\end{frame}

\begin{frame}
\frametitle{Порты и пакеты}
\begin{itemize}
  \item Порты используются для построения пакетов
  \item Зависимости и конфликты разрешаются pkg, а не make
  \item Стабильная ветка портов соответствует стабильной ветке пакетов
  \item Предлагать пользователям ставить ПО из пакетов
  \item Не запрещая установку из портов
\end{itemize} 
\end{frame}

\begin{frame}
\frametitle{Создание репозитория пакетов}
\begin{figure}[h!]
  \centering
  \includegraphics[width=0.95\textwidth]{q3.eps}
\end{figure}
\end{frame}

\begin{frame}
\frametitle{Установка пакетов в систему}
\begin{figure}[h!]
  \centering
  \includegraphics[width=0.95\textwidth]{q4.eps}
\end{figure}
\end{frame}

\begin{frame}
\frametitle{Проблемы в pkg}
\begin{itemize}
  \item Поддержка устаревших портов (например, без stage directory)
  \item Крайне ограниченная система зависимостей
  \item Примитивный алгоритм разрешения зависимостей (solver)
\end{itemize}
\end{frame}

\begin{frame}
\frametitle{Проблемы с алгоритмом solver'а}

\begin{itemize}
\item Нет поддержки конфликтов
\item Не работает с альтернативными пакетами
\item Может выполнять только одну задачу: установить, удалить или обновить
некоторые пакеты
\end{itemize}

\end{frame}


\begin{frame}
\frametitle{Существующие системы управления пакетами}

Другие системы управления пакетами используют различные подходы для решения
данной задачи:
\begin{itemize}
  \item \includegraphics[width=22pt]{suse.pdf} \hspace{5pt} \texttt{Zypper/SUSE}
  - построена на базе libsolv
  \item \includegraphics[width=22pt]{rpm.pdf} \hspace{5pt} \texttt{Yum/RedHat} -
  планирует использование libsolv
  \item \includegraphics[width=15pt]{debian.pdf} \hspace{12pt}
  \texttt{Apt/Debian} - базируестя на внутренней логике разрешения зависимостей
  и конфликтов
  \item \includegraphics[width=15pt]{arch.pdf} \hspace{12pt}
  \texttt{Pacman/Archlinux} - использует наивный алгоритм
\end{itemize}

\end{frame}

\begin{frame}[fragile]
\frametitle{Независимые solver'ы}
Для взаимодействия с внешними solver'ами, предполагается использовать формат
CUDF, разработанный в рамках исследовательского проекта Mancoosi
\url{http://mancoosi.org}:
\bigskip
{\small
	\begin{verbatim}
	package: devel/libblah
	version: 1
	depends: x11/libfoo

	package: security/blah
	version: 2
	depends: devel/libblah
	conflicts: security/blah-devel
	
	\end{verbatim}
}
\end{frame}

\begin{frame}
\frametitle{Но нам нужен собственный solver!}

Варианты:
\begin{itemize}
  \item Реализовать собственный алгоритм разрешения зависимостей?
  \pause
  \item Использовать существующие решения?
  \pause
  \item Использовать известный алгоритм?
  \pause
\end{itemize}
\bigskip
{\large SAT solver для управления пакетами}
\bigskip
\[
\overbrace{\underbrace{(x_1 \| \neg x_2 \|
x_3)}_{\text{\color{light-gray}Условие}} \& (x_3 \| \neg x_1) \&
(x_2)}^{\text{\color{light-gray}SAT выражение}}
\]
\end{frame}

\begin{frame}
\frametitle{Составление выражения SAT}
\begin{itemize}
  \item Assign a variable to each package: 
  package A $\to a_1$, package B $\to b_1$
  \item Interpret a request as a set of unary clauses:
  \begin{itemize}
    \item Install/Upgrade package A $\to (a_1)$
    \item Delete package B $\to (\neg b_1)$
  \end{itemize}
  \item Convert dependencies and conflicts to disjuncted clauses
\end{itemize}

\end{frame}

\begin{frame}
\frametitle{Converting dependencies and conflicts}
\begin{itemize}
  \item If package A depends on package B (versions $B_1$ and $B_2$), then we
  can either have package A not installed or any of B installed:
  \bigskip
\[
(\neg A \| B_1 \| B_2)
\]
\pause
  \item If we have a conflict between versions of B ($B_1$, $B_2$ and $B_3$)
 then we ensure that merely one version is installed:
  \bigskip
\[
\underbrace{(\neg B_1 \| \neg B_2) \& (\neg B_1 \| \neg B_3) \& (\neg B_2 \|
\neg B_3)}_{\text{\color{light-gray}Conflicts chain}}
\]
\end{itemize}
\end{frame}

\begin{frame}
\frametitle{The solving of SAT problem}

Some rules to follow to speed up SAT problem solving.
\begin{itemize}
  \item Trivial propagation - solve unary clauses
  \item Unit propagation - solve clauses with only a single unsolved variable
  \item Conflicts learning - if we assign some free variable and detect a
  conflict during unit propagation, we can fallback and learn that this variable
  must be negated
  \item Package specific assumptions.
\end{itemize}
\end{frame}

\begin{frame}
\frametitle{SAT problem propagation}
\begin{itemize}
  \item Trivial propagation - direct install or delete rules
  \bigskip
  \[
  (\neg A \| B) \& \underbrace{(A)}_{\textit{\color{light-gray}true}} \&
  \underbrace{(\neg C)}_{\textit{\color{light-gray}false}} \& (\neg A \| \neg D)
  \]
  \pause
  \item Unit propagation - simple depends and conflicts
  \bigskip
  \[
  \overbracket{\underbrace{(\neg A \| B)}_{B \rightarrow
  \textit{\color{light-gray}true}}}^{\text{\color{light-gray}Dependency}} \&
  \overbrace{(A)}^{\textit{\color{light-gray}true}} \& \overbrace{(\neg
  C)}^{\textit{\color{light-gray}false}} \&
  \overbracket{\underbrace{(\neg A \| \neg D)}_{D \rightarrow
  \textit{\color{light-gray}false}}}^{\text{\color{light-gray}Conflict}}
  \]
  \
\end{itemize}
\end{frame}

\begin{frame}
\frametitle{Conflicts driven learning}
To handle alternatives it is required to test all variables unassigned:
\begin{enumerate}
  \item full depth-first enumeration of possible values
  \item fallback if a conflict found
  \item remember which assignment caused conflict
  \item make negative assignment for the learned variable and go to the first
  step
\end{enumerate}
\end{frame}

\begin{frame}
\frametitle{Package specific assumptions}
Pure SAT solvers cannot deal with package management as they do not consider
several packages peculiarities:
\begin{itemize}
  \item try to keep installed packages (if no direct conflicts)
  \item do not install packages if they are not needed
  \item prefer high priority packages and repositories over low priority ones
\end{itemize}
These options also improve SAT performance providing a good initial assignment.
\end{frame}

\begin{frame}
\frametitle{Packages universe}
We convert all packages involved to a packages universe of the following
structure:
\begin{figure}[h!]
  \centering
  \includegraphics[height=0.7\textheight]{q5.eps}
\end{figure}
\end{frame}

\begin{frame}
\frametitle{Package management task}
\begin{itemize}
  \item A request is splitted to install/upgrade and delete requests which
  could be passed simultaneously to the solver
  \item A conflicts between packages are detected with a repository creation
  \item All depends, reverse and conflicts of the requested packages are
  analyzed and the package universe is created
  \item Each package is defined by its name and the digest of significant
  fields (version, options and so on)
\end{itemize}
\end{frame}

\begin{frame}
\frametitle{Solvers and Pkg}
\begin{itemize}
  \item Pkg may pass the formed universe to an external CUDF solver:
  \begin{itemize}
    \item convert versions
    \item format request
    \item parse output
  \end{itemize}
  \item Alternatively the internal SAT solver may be used:
  \begin{itemize}
    \item convert the universe to SAT problem
    \item formulate request
    \item ???
    \item PROFIT
  \end{itemize} 
\end{itemize}
\end{frame}

\begin{frame}
\frametitle{Perspectives}

\begin{itemize}
  \item Using pkg solver for ports management
  \item Better support of multiple repositories
  \item Test different solvers algorithms using CUDF
  \item New dependencies and conflicts format
  \item Provides and alternatives
\end{itemize}
\end{frame}

\begin{frame}
\frametitle{New dependencies format}
\large{\[libblah >= 1.0 +option_1, +option_2 \| libfoo != 1.1\]}
\begin{itemize}
  \item Can depend on normal packages and virtual packages (provides)
  \item Easy to define the concrete dependency versions
  \item Alternative dependencies
\end{itemize}
\begin{figure}[h!]
  \centering
  \includegraphics[height=0.3\textheight]{q6.eps}
\end{figure}
\end{frame}

\begin{frame}
\frametitle{Alternatives}
\begin{itemize}
  \item Used to organize packages with the same functionality (e.g.
  web-browser)
  \item May be used to implement virtual dependencies (provides/requires)
\end{itemize}
\begin{figure}[h!]
  \centering
  \includegraphics[height=0.4\textheight]{q7.eps}
\end{figure}
\end{frame}

\begin{frame}
\begin{center}
\includegraphics{logo.pdf}
{\Large Thank you for your attention!} \\
\emph{Questions?} \\[4pt]
\url{vsevolod@FreeBSD.org}
\end{center}
\end{frame}

\end{document}
